\documentclass{ximera}[font=12pt]


\title{Adding Integers: Same Signs}   % Use xmlatex all -s to compile when SERVE refuses to work.
\author{Kelly Stady}
\license{CC: 0}         % replace with an appropriate license, or set it in xmPreamble

\begin{document}

\begin{abstract}

\end{abstract}

\maketitle
\label{xim:AddIntSameSActivity}


\section*{Adding Integers with the Same Sign}

When adding numbers with the same sign, there is no ``battle.''  Instead, the numbers \textbf{join forces} to 
create a larger asset or a deeper debt.

\begin{example} \textcolor{glaucous}{\textbf{Joining Assets}}

If you have $\$400$ in your savings account ($+400$) and you win $\$150$ ($+150$), you have a net balance of \$550.
\[ 400 + 150 = 550 \]

\end{example}

\begin{example} \textcolor{glaucous}{\textbf{Joining Debts}}

If you owe your brother $\$10$ ($-10$) and then borrow another $\$25$ ($-25$), your net balance is a debt of $\$35.$  
So, you are falling deeper into debt.
\[ -10 + (-25) = -35 \]

\end{example}


We can now write a formal mathematical rule for adding integers with the same sign.  Once again, we will use absolute value.


\begin{procedure}[\textbf{Steps for Adding Integers with the Same Signs}]

\begin{enumerate} 
    \item Find the \textbf{absolute value} or \textbf{strength} of each number.

     \item \textbf{Add} the absolute values.

     \item Write the answer with the common sign.
\end{enumerate}
\end{procedure}


\begin{example} 
    
    Use the three-step process to find the sum $-50 + (-30).$

\begin{explanation}
    
    \begin{description}[style=multiline, leftmargin=5em, font=\bfseries]
    
    \item[Step 1] \textbf{Find the Absolute Values (Strengths)}: \\
    Fnd the strength of each number.
    \[ |-50| = 50 \quad \text{and} \quad |30| = 30 \]

    \item[Step 2] \textbf{Add the Absolute Values (Join Forces)}: \\
    Since the signs are the same, we add their absolute values.
    \[ 50 + 30 = 80 \]

    \item[Step 3] \textbf{Determine the Sign and Write the Answer}: \\
    Since both numbers are negative, the answer is negative.
    \[ -50 + (-30) = -80 \]

\end{description}
\end{explanation}
\end{example}


\begin{problem}
Apply the steps to find the sum $-15 + (-40).$

~ \\[1em]

\begin{description}[style=multiline, leftmargin=5em, font=\bfseries]
    
    \item[Step 1] \textbf{Find the Absolute Values (Strengths)}: \\
    
    \[ |-15| = \answer{15} \quad \text{and} \quad |-40| = \answer{40} \]

    \item[Step 2] \textbf{Add the Absolute Values (Join Forces)}: \\
    
    \[ 15 + 40 = \answer{55} \]

    \item[Step 3] \textbf{Determine the Sign and Write the Answer}: \\
    
     \ $-15 + (-40) = \answer{-55}$

\end{description}
\end{problem}


\begin{problem}
A research submarine is exploring the ocean floor.  It starts at a depth of $5$ meters below sea level.  The pilot then 
navigates the submarine down another $8$ meters.  Find the submarine's final depth relative to sea level.

\[ -5 + (-13) = \answer{-18} \ \textrm{meters} \]

\end{problem}


\begin{problem}
During a football game, the home team's running back gained $6$ yards on the first play of the drive. 
On the very next play, he charged forward an additional $7$ yards.  Find the total number of yards 
gained by the running back over these two plays.

\[ 6 + 7 = \answer{13} \ \textrm{yards} \]

\end{problem}


\begin{problem}
On a cold winter night in Cleveland, the starting temperature at midnight was $-1^\circ\text{F}$.  Over the next three 
hours, another cold front moved through the area and the temperature dropped by $9^\circ\text{F}.$  Find the 
temperature at 3am.

\[ -1 + (-9) = \answer{-10} \ ^\circ\text{F} \]
\end{problem}



\end{document}
